% Preamble for collective_agents.tex
\usepackage[T1]{fontenc}
\usepackage{lmodern}
\usepackage[letterpaper,margin=1in]{geometry}
\usepackage{microtype}
% A third line-breaking pass for paragraphs that hold unbreakable material (script names, long inline formulas).
\setlength{\emergencystretch}{2em}
\usepackage{amsmath,amssymb,amsthm,mathtools}
\usepackage{aliascnt}
\usepackage{bm}
\usepackage{booktabs}
\usepackage{array}
% Tables longer than one page (the notation and summary tables) break across pages.
\usepackage{longtable}
\setlength{\LTcapwidth}{\textwidth}
\usepackage{xcolor}
\usepackage{tikz}
\usetikzlibrary{arrows.meta,positioning,calc,fit,backgrounds,shapes.geometric}
\usepackage{pgfplots}
\pgfplotsset{compat=1.18}
\usepackage[authoryear,round]{natbib}
\usepackage[colorlinks=true,linkcolor=blue!50!black,citecolor=green!35!black,urlcolor=blue!50!black]{hyperref}
\usepackage[capitalize,nameinlink]{cleveref}

% Theorem-like environments share the theorem counter through aliased counters,
% so that cleveref names each one correctly (it names shared counters by the
% counter otherwise).
\theoremstyle{plain}
\newtheorem{theorem}{Theorem}[section]
\newcommand{\sharedtheorem}[2]{\newaliascnt{#1}{theorem}\newtheorem{#1}[#1]{#2}\aliascntresetthe{#1}}
\sharedtheorem{proposition}{Proposition}
\sharedtheorem{lemma}{Lemma}
\sharedtheorem{corollary}{Corollary}
\sharedtheorem{conjecture}{Conjecture}
\theoremstyle{definition}
\sharedtheorem{definition}{Definition}
\sharedtheorem{assumption}{Assumption}
\sharedtheorem{example}{Example}
\sharedtheorem{observation}{Numerical observation}
\sharedtheorem{openproblem}{Open problem}
\newtheorem{hypothesis}{Hypothesis}
\renewcommand{\thehypothesis}{H\arabic{hypothesis}}
\newtheorem{premise}{Assumption}
\renewcommand{\thepremise}{A\arabic{premise}}
\newtheorem{numcheck}{Check}[section]
\theoremstyle{remark}
\sharedtheorem{remark}{Remark}

\crefname{theorem}{Theorem}{Theorems}
\crefname{proposition}{Proposition}{Propositions}
\crefname{lemma}{Lemma}{Lemmas}
\crefname{corollary}{Corollary}{Corollaries}
\crefname{conjecture}{Conjecture}{Conjectures}
\crefname{definition}{Definition}{Definitions}
\crefname{assumption}{Assumption}{Assumptions}
\crefname{example}{Example}{Examples}
\crefname{observation}{Numerical observation}{Numerical observations}
\crefname{openproblem}{Open problem}{Open problems}
\crefname{remark}{Remark}{Remarks}
\crefname{hypothesis}{Hypothesis}{Hypotheses}
\crefname{premise}{Assumption}{Assumptions}
\crefname{numcheck}{Check}{Checks}

% Notation
\newcommand{\E}{\mathbb{E}}
\newcommand{\R}{\mathbb{R}}
\newcommand{\Z}{\mathbb{Z}}
\newcommand{\Ncal}{\mathcal{N}}
\newcommand{\Cov}{\operatorname{Cov}}
\newcommand{\Var}{\operatorname{Var}}
\newcommand{\KL}{D_{\mathrm{KL}}}
\newcommand{\tr}{\operatorname{tr}}
\newcommand{\rank}{\operatorname{rank}}
\newcommand{\diag}{\operatorname{diag}}
\newcommand{\GL}{\mathrm{GL}}
\newcommand{\arccosh}{\operatorname{arccosh}}
\newcommand{\arcsinh}{\operatorname{arcsinh}}
\newcommand{\Corr}{\operatorname{Corr}}
\newcommand{\Hol}{\operatorname{Hol}}
\newcommand{\Gibbs}{\operatorname{Gibbs}}
\newcommand{\perc}{\operatorname{perc}}
\newcommand{\cut}{\operatorname{cut}}
\newcommand{\one}{\mathbf{1}}
\newcommand{\Kg}{\mathcal{K}}
\newcommand{\sI}{\mathsf{I}}
\newcommand{\sS}{\mathsf{S}}
\newcommand{\sA}{\mathsf{A}}
\newcommand{\sE}{\mathsf{E}}
\newcommand{\sB}{\mathsf{B}}
\newcommand{\xin}{\varkappa_{\mathrm{in}}}
\newcommand{\xout}{\varkappa_{\mathrm{out}}}
\newcommand{\Agency}{\mathfrak{A}}
\newcommand{\Lstar}{L^{\ast}}
